\chapter{Outlook: three doors and a costed path}\label{ch:outlook}

The preceding chapters were written in the order the work was done, which means
they read as a narrowing. That narrowing is the result. This chapter states it
positively: after Chapter~\ref{ch:closure}, everything that could possibly
improve a Goldbach-desert record passes through one of exactly three doors, and
each door has a maximal form, a price, and a verdict.

\section{The closure argument}

Why only three? Because the three ways to change $e^{E}$ are to change the
mathematics that sets $E$, to change the rate at which candidates are examined,
or to leave the paradigm in which $E$ is defined. Chapter~\ref{ch:frontier}
showed the first is a rounding gap, Chapter~\ref{ch:engineering} measured the
second, and Chapter~\ref{ch:closure} closed the third within a concrete
language. Two auxiliary facts complete the argument, and both are worth stating
because they eliminate whole families of otherwise-attractive ideas:

\begin{itemize}
\item \emph{Divisor information is fully harvested.} Any congruence cleverness
applied to the search parameter $k$ merely re-labels sub-progressions that the
sieve already scores. There is no free information on the $k$ side.
\item \emph{Correlation is covering} (Proposition~\ref{prop:shared}), and
beyond divisors the P5 measurement caps all remaining structure at $0.1$ nats.
So ``luck-shaping'' schemes have no purchase.
\end{itemize}

\section{Door 1: the rounding gap (mathematics)}

The exponent is entirely an integrality gap (Corollary~\ref{cor:rounding}), so
better rounding is worth records. Measured prize: greedy's $|U|=867$ against an
annealed ensemble threshold near $758$, i.e.\ $0$ to $2.6$ nats---a factor of up
to $14$, compounding into \emph{every} future rung.

Verdict from Chapter~\ref{ch:frontier}: resistant. Exact windows (LP-certified),
half-cover exact blocks, and BP decimation all leave $867$ untouched, and the
landscape is a field of isolated basins in which the annealed threshold is
probably not achievable. What has \emph{not} been tried is a different
\emph{ensemble} rather than a better search of the same one. Specifically:

\begin{enumerate}
\item \textbf{Banded construction with an exact endgame.} Solve the small moduli
exactly, choose the middle band with second-moment control on pairwise overlaps,
and finish the large moduli as an exact weighted matching against the
currently-uncovered set. This imitates \emph{why} greedy wins---sequential
adaptivity---while removing its myopia at the point where overlap waste is
greatest.
\item \textbf{Import Erd\H{o}s--Rankin technology constructively.} The
FGKMT machinery~\cite{FGKMT} is, at bottom, a method for rounding a fractional
cover of an interval with quantified loss; we used it in
Theorem~\ref{thm:lower} for a lower bound but never as a construction. Its loss
is $o(1)$ in an asymptotic regime that is not ours (their moduli reach $y/2$;
ours stop at $457$), so the question ``what residual fraction is achievable at
$\pi(Q)\approx10^{4}$ with a budget of $\theta(457)$?'' is genuinely open and
thesis-sized.
\item \textbf{One-step replica-symmetry-breaking.} Needed to get an honest
quenched threshold; naive belief propagation demonstrably slides toward the
fractional fantasy (\S\ref{sec:door1}).
\item \textbf{Post the instance as a public optimisation challenge.} The
benchmark is small, self-contained, and machine-checkable. Problems of exactly
this shape attract strong heuristic solvers when posed publicly.
\end{enumerate}

\section{Door 2: throughput (compute)}

Records are logarithmic in compute: each $10\times$ buys $2.3$ nats, which is
about four digits on the threshold game. Door 2 is therefore unglamorous and
completely reliable. Three forms, priced.

\paragraph{Engineering already in hand.} Porting the selective-testing scheme of
\S\ref{sec:engine} to the GPU fleet is worth $2$--$3\times$; batch
remainder-tree sieving~\cite{Bernstein} to depth $10^{9}$--$10^{10}$ is worth
perhaps $1.5\times$; replacing the fixed skip fraction with the exact marginal
stopping rule adds $5$--$15\%$. None of this is research.

\paragraph{Volunteer computing.} The search is a progress-free Poisson lottery
with millisecond verification and no shared state---structurally the best-shaped
volunteer workload that exists, and the exact shape on which
PrimeGrid~\cite{PrimeGrid} and its relatives hold essentially every record in
the neighbouring Sierpi\'nski--Riesel problems, which are themselves covering
problems. For calibration: a peak-era Folding@home-scale network is roughly
$2\times10^{5}$ times our two-GPU fleet, which would put $Q=200\,000$
(Table~\ref{tab:ladder}) about \emph{two days} away; even $1\%$ of that scale
delivers it in months.

\paragraph{Silicon.} The workload is $95\%$ one instruction---a
$660$-bit Montgomery modular exponentiation with a fixed base. A
crypto-mining-style ASIC plausibly buys $10$--$100\times$ an A100 per die at
much better energy per test, so a single rack is worth $e^{5.5}$ to $e^{8}$ over
today's fleet; an FPGA farm is the no-NRE middle step. This is the
capital-for-nats trade, and it is priced in the single-digit millions.

\section{Door 3: the paradigm}

Closed, within the language $L$ (Chapter~\ref{ch:closure}). The deliverable here
has therefore flipped from \emph{find a mechanism} to \emph{finish the theorem
and keep trying to break it}: complete the rigour of Lemma~\ref{lem:entropy},
settle (H2) or publicise it, and keep the falsification engine running as new
schema families are proposed. If the closure claim is ever refuted, the payoff
is enormous and immediate---but the sensible planning assumption is that it will
not be.

\section{Do we need a quantum computer?}

Grover search~\cite{Grover} halves the exponent: $e^{E}$ classical candidates
become $e^{E/2}$ coherent oracle calls. Since $E=37.7$ at $Q=200\,000$ and
$44.1$ for sub-$100$, this looks decisive. It is not, and the arithmetic is
worth doing once so that it need not be revisited.

One oracle call must evaluate the whole predicate reversibly: roughly $|U|$
modular exponentiations of $660$-bit numbers with \emph{no early exit}, since
the classical trick of stopping at the first passing complement is unavailable
in superposition. At fault-tolerant resource-estimate
scales~\cite{GidneyEkera} this is on the order of $10^{5}$--$10^{6}$ seconds per
call. A classical candidate costs $1.25\times10^{-6}$ seconds on our fleet.
Writing $C_q,W_q$ and $C_c,W_c$ for the per-operation costs and the widths
(numbers of machines), quantum wins only when
\[
E\;>\;E^{*}\;=\;2\log\Bigl[\frac{C_q}{W_q}\cdot\frac{W_c}{C_c}\Bigr].
\]

\begin{table}[h]
\centering\small
\caption{The Grover crossover. Compare against $E=37.7$ ($Q=200\,000$) and
$E=44.1$ (sub-$100$).}\label{tab:quantum}
\begin{tabular}{lr}
\toprule
Assumption & $E^{*}$\\
\midrule
One fault-tolerant machine, $C_q=10^{6}$\,s & $\approx55$\\
One thousand such machines & $\approx41$\\
One thousand machines \emph{and} a $1000\times$ cheaper oracle & $\approx27$\\
\bottomrule
\end{tabular}
\end{table}

With any non-miraculous constants the crossover sits at or above the hardest
targets we care about---and at the crossover both approaches are equally
infeasible. The honest answer is therefore: \emph{no}. You would need thousands
of much better quantum computers than anyone projects merely to \emph{tie} with
GPUs on this problem. Quantum walks for subset-sum face the same arithmetic in
Chapter~\ref{ch:barriers}'s setting. We park the idea with numbers attached
rather than dismissing it, so that it can be revisited if $C_q$ ever approaches
seconds.

\section{Operation Staircase: a costed path to \texorpdfstring{$\g>200\,000$}{g > 200000}}

Putting the doors together gives a concrete plan. Its shape is dictated by
Table~\ref{tab:ladder}: rungs are geometrically spaced, so climbing all of them
costs about $4\%$ more than jumping to the top, while banking a certified record
and a fresh calibration at every step.

\begin{enumerate}
\item \textbf{Now.} $Q=140\,009$ ($E\approx25$) on the existing fleet: days.
Bank it.
\item \textbf{Plus one week.} Land the engine work (selective testing on GPU,
marginal stopping rule): $2$--$3\times$. Then $Q=150\,000$ ($E=27.3$): about
four fleet-days. Bank it.
\item \textbf{Plus one month.} $Q\approx165\,000$ ($E\approx30$): two to three
fleet-weeks with the upgraded engine. In parallel, Door 1's banded construction
and the Erd\H{o}s--Rankin import run as theory work; any $\ge1$ nat result
compounds into every later rung.
\item \textbf{Plus one quarter.} $Q=175\,000$ ($E=32.1$): $1.4$ fleet-years raw,
so two to four months with the engine work and modest fleet growth. Decision
gate: if Door 1 has delivered $\ge2$ nats, stay in-house; otherwise open Door 2
at scale.
\item \textbf{The wall.} $Q=190\,000$--$200\,000$ ($E=35.4$--$37.7$) needs about
$e^{5.5}$ beyond step 4. Three interchangeable closers: a few hundred swarm or
volunteer GPUs for a quarter; an ASIC/FPGA rack; or four nats of Door-1
mathematics plus a $10\times$ fleet. Any one suffices; any two make it
comfortable.
\end{enumerate}

Expected profile: records at $140$k and $150$k within weeks, $165$k within a
month or two, $175$k inside a quarter, and $200\,000$ in two to four
quarters---with the long pole attackable by money, community, or mathematics,
whichever the programme can actually raise. Throughout, every rung is
dual-verified by the pipeline of \S\ref{sec:certification}, and parallel workers
claim disjoint variant sets in a committed ledger. That last point is not
bureaucratic: one of our own campaigns was stopped mid-flight for duplicating
another effort's target, which is a pure waste that coordination prevents.

\section{Open problems}

We close with the problems this thesis leaves in a state where someone else can
pick them up. The first two would change the subject; the rest would advance it.

\begin{enumerate}
\item \textbf{Sub-modexp compositeness detection} (Conjecture~\ref{conj:h2}).
Is there a randomised test with constant soundness on rough inputs, cheaper than
one modular exponentiation? Or a lower bound in any realistic model? A positive
answer speeds up every prime-hunting project in existence; a negative one
completes Lemma~\ref{lem:entropy}'s third leg.

\item \textbf{The rounding gap for prime-residue covers.} Given $\pi(Q)\approx
10^{4}$ prime offsets and moduli up to $457$, what residual fraction is
achievable? Greedy gives $7.56\%$, the fractional optimum $0\%$, and the annealed
ensemble threshold about $6.6\%$. Where is the quenched truth, and is it
constructive?

\item \textbf{CRT list-recovery beyond the Johnson radius} for structured,
solution-abundant instances (Chapter~\ref{ch:barriers}). Our experiments say the
radius is a sharp algorithmic boundary in practice; a proof or an algorithm
would both be significant.

\item \textbf{Modular subset-sum with per-position choices at density
$\approx1$}, where $2^{82}$ solutions exist among $2^{3893}$ designs.

\item \textbf{Mechanisms outside $L$}: certify $\N-q$ composite by something
that is neither a planted divisor identity nor per-instance divisor discovery.
This is the honest frontier of Chapter~\ref{ch:closure}'s claim.

\item \textbf{The smallest prescribed desert.} $S(q)$, the least even $\N$ with
$\g(\N)=q$~\cite{GLvdtR}, packages into $T(Q)=\min_{q>Q}S(q)$; no construction
here determines an exact minimum, and exhaustive verification stops at
$4\cdot10^{18}$ while our constructions start at $10^{150}$. The intervening
$130$ orders of magnitude are entirely unexplored.

\item \textbf{Closing the asymptotic sandwich.} Between
$\log\N\log_2\N\log_4\N/\log_3\N$ (rigorous, Theorem~\ref{thm:lower}) and
$(\log\N)^{2}\log\log\N$ (conjectural~\eqref{eq:GLR}) there is a factor of about
$\log\N$. Does covering only prime offsets---rather than a whole
interval---asymptotically beat the prime-gap scale?
\end{enumerate}

\section{A closing note on method}

This thesis was produced with heavy delegation to automated agents: search
campaigns, cover optimisation, literature checking, and an adversarial
falsification engine. Three practices did more than any individual result to
keep the output trustworthy, and we recommend them without reservation.

First, \emph{calibrate before believing}. The failure exponent was checked
against $1.2\times10^{9}$ candidates before it was used to make claims; the
belief-propagation implementation was validated against exhaustive enumeration
on a small instance before it was run on a large one. Both checks later caught
real errors.

Second, \emph{make negative results first-class}. Three campaigns produced no
record. Reported with their survival probabilities they became measurements
rather than embarrassments, and they now carry a substantial share of the
model's empirical support.

Third, \emph{let the adversary audit the referee}. The single most useful event
in Chapter~\ref{ch:closure} was an evolutionary search discovering a $50\times$
error in our own scorer. The audit protocol that caught it---candidates emit
data, never scores; grading in a separate trusted process; every quantity
re-derived from primitives; a ``pass'' is hand-audited---is cheap, general, and
worth more than the theorem it was built to defend. We also retracted two claims
in the course of this work (\S\ref{sec:retraction}, Remark~\ref{rem:units}) and
recorded both in place rather than quietly deleting them, which is the same
discipline pointed at ourselves.
